\section{Robustness and Limitations}
\label{sec:robustness}
Our findings are subject to several important limitations that warrant disclosure:
\begin{enumerate}
    \item \textbf{Evaluation Window Size:} The computationally intensive nature of the Optuna-tuned walk-forward Neural-HAR model restricted its out-of-sample evaluation to approximately 275-287 trading days. This window represents a recent, relatively calm market environment and excludes the extreme volatility clustering of the 2020 COVID-19 pandemic. The baseline HAR model's headline statistics were computed over a much longer multi-year sample; our direct comparative tests (DM and MCS) align the samples to the shorter window, but the robustness of the neural model across rare macro shocks remains untested in this paper.
    \item \textbf{Regime Sample Sizes:} The regime-conditional analysis partitions the data into subsets, resulting in relatively small sample sizes per regime (e.g., 155 to 260 observations for SPX and NIFTY low/medium states). This limits the statistical power of the regression coefficient tests in Table \ref{tab:regime_coefs}.
    \item \textbf{Implementation Costs:} The economic significance tests reveal that the predictive edge of sentiment and deep learning models is highly sensitive to transaction costs. At 5 bps slippage, the alpha generated by the Neural-HAR model is fully consumed by its turnover, limiting its utility for yield-seeking portfolios, though it remains viable for downside risk control.
\end{enumerate}
